\documentclass[11pt,a4paper]{article}

% ============================================================
%                         PACKAGES
% ============================================================

\usepackage{amsmath,amssymb,mathtools}
\usepackage{tikz}
\usepackage{tikz-3dplot}
\usetikzlibrary{arrows.meta}

\usepackage{geometry}
\usepackage{xcolor}
\usepackage{tcolorbox}
\usepackage{pagecolor}
\usepackage{fontspec}
\usepackage{microtype}
\usepackage{hyperref}

% ============================================================
%                         CONFIGURATION
% ============================================================

\geometry{
    top=2.2cm,
    bottom=2.2cm,
    left=2.3cm,
    right=2.3cm
}

% ============================================================
%                         COULEURS
% ============================================================

\definecolor{background}{HTML}{010402}
\definecolor{text}{HTML}{F0F0F0}
\definecolor{muted}{HTML}{999999}

\definecolor{accent}{HTML}{64B5F6}
\definecolor{accent2}{HTML}{BA68C8}
\definecolor{green}{HTML}{81C784}
\definecolor{orange}{HTML}{FFB74D}
\definecolor{red}{HTML}{E57373}

\definecolor{boxbg}{HTML}{050505}
\definecolor{grid}{HTML}{080B09}

\pagecolor{background}
\color{text}

% ============================================================
%                         POLICE
% ============================================================

\setmainfont{Latin Modern Roman}

% ============================================================
%                         HYPERLIENS
% ============================================================

\hypersetup{
    colorlinks=true,
    linkcolor=accent,
    urlcolor=accent,
    citecolor=accent
}

% ============================================================
%                         COMMANDES
% ============================================================

\newcommand{\R}{\mathbb{R}}

\newcommand{\vect}[1]{%
    \begin{pmatrix}
        #1
    \end{pmatrix}%
}

\newcommand{\sectiontitle}[1]{\section{#1}}

% ============================================================
%                         BOXES
% ============================================================

\newtcolorbox{definitionbox}{%
    colback=boxbg,
    colframe=accent!70!black,
    coltext=text,
    boxrule=0.7pt,
    arc=3mm,
    left=5mm,
    right=5mm,
    top=3mm,
    bottom=3mm
}

\newtcolorbox{importantbox}{%
    colback=boxbg,
    colframe=orange!70!black,
    coltext=text,
    boxrule=0.7pt,
    arc=3mm,
    left=5mm,
    right=5mm,
    top=3mm,
    bottom=3mm
}

% ============================================================
%                         DOCUMENT
% ============================================================

\begin{document}

% ============================================================
%                         PAGE DE GARDE
% ============================================================

\begin{center}

\vspace*{2cm}

{\fontsize{38}{45}\selectfont\bfseries Matrices 3D\par}

\vspace{0.45cm}

{\fontsize{20}{25}\selectfont Géométrie et transformations dans l'espace\par}

\vspace{1cm}

\tdplotsetmaincoords{68}{122}

\begin{tikzpicture}[tdplot_main_coords,scale=1.6]

    % Axes
    \draw[-{Latex[length=2.5mm]},line width=0.8pt,color=muted]
        (0,0,0)--(2.1,0,0) node[right] {$x$};

    \draw[-{Latex[length=2.5mm]},line width=0.8pt,color=muted]
        (0,0,0)--(0,1.9,0) node[above] {$y$};

    \draw[-{Latex[length=2.5mm]},line width=0.8pt,color=muted]
        (0,0,0)--(0,0,2.0) node[above] {$z$};

    % Cube
    \coordinate (A) at (0.15,0.15,0.1);
    \coordinate (B) at (1.15,0.15,0.1);
    \coordinate (C) at (1.15,1.15,0.1);
    \coordinate (D) at (0.15,1.15,0.1);
    \coordinate (E) at (0.15,0.15,1.1);
    \coordinate (F) at (1.15,0.15,1.1);
    \coordinate (G) at (1.15,1.15,1.1);
    \coordinate (H) at (0.15,1.15,1.1);

    \draw[color=accent!35,line width=0.85pt]
        (A)--(B)--(C)--(D)--cycle;
    \draw[color=accent!55,line width=0.85pt]
        (E)--(F)--(G)--(H)--cycle;
    \draw[color=accent!45,line width=0.8pt]
        (A)--(E) (B)--(F) (C)--(G) (D)--(H);

    \draw[-{Latex[length=2.6mm]},color=orange,line width=1.35pt]
        (A)--(G);

    \fill[color=orange] (G) circle (1.7pt);

    \node[color=orange,anchor=west,inner sep=3pt]
        at (1.20,1.17,1.12) {$\vec v$};

\end{tikzpicture}

\vfill

{\color{muted}De la géométrie spatiale aux transformations linéaires\par}

\vspace{1cm}

\end{center}

\newpage

% ============================================================
%                         INTRODUCTION
% ============================================================

\sectiontitle{Introduction}

En dimension 2, une matrice $2\times2$ transforme les coordonnées
d'un vecteur du plan.

En dimension 3, la même idée s'étend naturellement à trois coordonnées :

\[
\vec v=
\begin{pmatrix}
 x\\
 y\\
 z
\end{pmatrix}.
\]

\begin{definitionbox}
Une matrice $3\times3$ représente une transformation linéaire
de $\mathbb{R}^3$ vers lui-même.
\[
\boxed{
\mathbb{R}^3
\xrightarrow{\;A\;}
\mathbb{R}^3
}
\]
\end{definitionbox}

On va partir du repère spatial, construire les rotations autour
des trois axes, puis voir comment les matrices permettent de
transformer des objets 3D.

% ============================================================
%                         REPERE
% ============================================================

\sectiontitle{Le repère dans l'espace}

Dans le plan, deux axes suffisent.

Dans l'espace, on ajoute l'axe $z$.

Un point est donc décrit par :

\[
P=(x,y,z).
\]

\begin{center}

\tdplotsetmaincoords{70}{120}

\begin{tikzpicture}[tdplot_main_coords,scale=2.45]

    % Axes
    \draw[-{Latex[length=2.5mm]},line width=0.8pt,color=muted]
        (-0.15,0,0)--(1.8,0,0) node[right] {$x$};
    \draw[-{Latex[length=2.5mm]},line width=0.8pt,color=muted]
        (0,-0.15,0)--(0,1.7,0) node[above] {$y$};
    \draw[-{Latex[length=2.5mm]},line width=0.8pt,color=muted]
        (0,0,-0.15)--(0,0,1.8) node[above] {$z$};

    % Plan xy
    \draw[color=grid!95]
        (0,0,0)--(1.3,0,0)--(1.3,1.3,0)--(0,1.3,0)--cycle;

    \draw[color=grid!75]
        (0.43,0,0)--(0.43,1.3,0)
        (0.86,0,0)--(0.86,1.3,0)
        (0,0.43,0)--(1.3,0.43,0)
        (0,0.86,0)--(1.3,0.86,0);

    % Point et vecteur
    \coordinate (P) at (1.05,0.72,1.0);
    \coordinate (Q) at (1.05,0.72,0);

    \draw[-{Latex[length=2.7mm]},color=accent,line width=1.35pt]
        (0,0,0)--(P);

    \draw[dashed,color=muted!75,line width=0.7pt]
        (P)--(Q)
        (Q)--(1.05,0,0)
        (Q)--(0,0.72,0);

    \fill[color=orange] (P) circle (1.9pt);

    \node[color=orange,anchor=west,inner sep=3pt]
        at (1.1,0.77,1.03) {$P=(x,y,z)$};

    \node[color=accent,anchor=west,inner sep=2pt]
        at (0.5,0.04,0) {$x$};

    \node[color=accent2,anchor=west,inner sep=2pt]
        at (1.08,0.38,0) {$y$};

    \node[color=green,anchor=west,inner sep=2pt]
        at (0.05,0.04,0.7) {$z$};

\end{tikzpicture}

\end{center}

% ============================================================
%                         NORME
% ============================================================

\sectiontitle{La longueur d'un vecteur 3D}

Pour

\[
\vec v=
\begin{pmatrix}
 x\\
 y\\
 z
\end{pmatrix},
\]

la norme euclidienne est :

\[
\boxed{
\|\vec v\|
=
\sqrt{x^2+y^2+z^2}
}
\]

C'est l'extension tridimensionnelle du théorème de Pythagore.

% ============================================================
%                         SPHERE
% ============================================================

\newpage

\sectiontitle{La sphère unité et les angles 3D}

En 2D, le cercle unité vérifie :

\[
x^2+y^2=1.
\]

En 3D, la sphère unité vérifie :

\[
\boxed{x^2+y^2+z^2=1}
\]

On peut décrire un point de la sphère avec deux angles :

\[
\varphi=\text{angle autour de l'axe }z,
\]

et

\[
\theta=\text{angle mesuré depuis l'axe }z.
\]

Une paramétrisation classique est :

\[
\boxed{
\begin{aligned}
x&=\sin\theta\cos\varphi,\\
y&=\sin\theta\sin\varphi,\\
z&=\cos\theta.
\end{aligned}}
\]

\begin{center}

\tdplotsetmaincoords{66}{122}

\begin{tikzpicture}[tdplot_main_coords,scale=2.4]

    \def\th{55}
    \def\ph{38}

    % Axes
    \draw[-{Latex[length=2.5mm]},line width=0.8pt,color=muted]
        (-1.25,0,0)--(1.4,0,0) node[right] {$x$};
    \draw[-{Latex[length=2.5mm]},line width=0.8pt,color=muted]
        (0,-1.25,0)--(0,1.4,0) node[above] {$y$};
    \draw[-{Latex[length=2.5mm]},line width=0.8pt,color=muted]
        (0,0,-0.15)--(0,0,1.4) node[above] {$z$};

    % Equateur
    \draw[color=accent!30,line width=0.75pt]
        plot[domain=0:360,samples=100]
        ({cos(\x)},{sin(\x)},0);

    % Méridien
    \draw[color=accent2!28,line width=0.75pt]
        plot[domain=-90:90,samples=70]
        ({cos(\x)},0,{sin(\x)});

    % Parallèle
    \pgfmathsetmacro{\rp}{sin(\th)}
    \pgfmathsetmacro{\zp}{cos(\th)}

    \draw[color=green!35,line width=0.75pt]
        plot[domain=0:360,samples=90]
        ({\rp*cos(\x)},{\rp*sin(\x)},{\zp});

    % Point
    \pgfmathsetmacro{\px}{sin(\th)*cos(\ph)}
    \pgfmathsetmacro{\py}{sin(\th)*sin(\ph)}
    \pgfmathsetmacro{\pz}{cos(\th)}

    \coordinate (P) at (\px,\py,\pz);
    \coordinate (Q) at (\px,\py,0);

    % Rayon OP
    \draw[-{Latex[length=2.7mm]},color=accent,line width=1.35pt]
        (0,0,0)--(P);

    % Projection
    \draw[dashed,color=muted!75,line width=0.7pt]
        (P)--(Q);

    \draw[-{Latex[length=2.5mm]},color=accent2,line width=1.05pt]
        (0,0,0)--(Q);

    % Angle phi dans le plan xy
    \draw[color=green,line width=1.1pt]
        plot[domain=0:\ph,samples=35]
        ({0.34*cos(\x)},{0.34*sin(\x)},0);

    % Angle theta dans le plan vertical contenant OP
    \draw[color=orange,line width=1.1pt]
        plot[domain=0:\th,samples=40]
        ({0.38*sin(\x)*cos(\ph)},
         {0.38*sin(\x)*sin(\ph)},
         {0.38*cos(\x)});

    % Points
    \fill[color=orange] (P) circle (2pt);
    \fill[color=accent2] (Q) circle (1.6pt);

    % Labels
    \node[color=orange,anchor=west,inner sep=3pt]
        at (\px+0.06,\py+0.05,\pz+0.02) {$P$};

    \node[color=accent2,anchor=west,inner sep=3pt]
        at (\px,\py,0.04) {$P_{xy}$};

    \node[color=green,anchor=west,inner sep=2pt]
        at (0.36,0.12,0.02) {$\varphi$};

    \node[color=orange,anchor=west,inner sep=2pt]
        at (0.14,0.18,0.24) {$\theta$};

\end{tikzpicture}

\end{center}

Les deux angles sont représentés dans deux plans différents.

% ============================================================
%                         MATRICE 3x3
% ============================================================

\sectiontitle{Construire une matrice trois par trois}

Une matrice générale de dimension 3 est :

\[
\boxed{
A=
\begin{pmatrix}
 a&b&c\\
 d&e&f\\
 g&h&i
\end{pmatrix}}
\]

Appliquée à :

\[
\vec v=
\begin{pmatrix}
 x\\
 y\\
 z
\end{pmatrix},
\]

elle donne :

\[
\begin{pmatrix}
x'\\y'\\z'\end{pmatrix}
=
A
\begin{pmatrix}x\\y\\z\end{pmatrix}.
\]

Donc :

\[
\boxed{
\begin{aligned}
x'&=ax+by+cz,\\
y'&=dx+ey+fz,\\
z'&=gx+hy+iz.
\end{aligned}}
\]

\begin{importantbox}
Chaque coordonnée finale est une combinaison linéaire des trois
coordonnées initiales.
\end{importantbox}

% ============================================================
%                         ROTATION Z
% ============================================================

\sectiontitle{Rotation autour de l'axe z}

La rotation autour de $z$ conserve la coordonnée $z$ :

\[
\boxed{
R_z(\theta)=
\begin{pmatrix}
\cos\theta&-\sin\theta&0\\
\sin\theta&\cos\theta&0\\
0&0&1
\end{pmatrix}}
\]

\begin{center}

\tdplotsetmaincoords{68}{122}

\begin{tikzpicture}[tdplot_main_coords,scale=2.45]

    \def\a{58}
    \def\h{1.0}

    % Axes
    \draw[-{Latex[length=2.5mm]},line width=0.8pt,color=muted]
        (-1.25,0,0)--(1.5,0,0) node[right] {$x$};
    \draw[-{Latex[length=2.5mm]},line width=0.8pt,color=muted]
        (0,-1.25,0)--(0,1.5,0) node[above] {$y$};
    \draw[-{Latex[length=2.5mm]},line width=0.8pt,color=muted]
        (0,0,-0.1)--(0,0,1.45) node[above] {$z$};

    % Axe de rotation
    \draw[color=green,line width=1.1pt]
        (0,0,0)--(0,0,1.3);

    % Cercle de rotation
    \draw[color=accent2!30,line width=0.75pt]
        plot[domain=0:360,samples=90]
        ({cos(\x)},{sin(\x)},0);

    % Vecteurs
    \draw[-{Latex[length=2.7mm]},color=muted,line width=1.3pt]
        (0,0,0)--(1,0,\h);

    \draw[-{Latex[length=2.7mm]},color=accent,line width=1.35pt]
        (0,0,0)--({cos(\a)},{sin(\a)},\h);

    % Projections
    \draw[dashed,color=muted!75,line width=0.7pt]
        (1,0,0)--(1,0,\h);
    \draw[dashed,color=muted!75,line width=0.7pt]
        ({cos(\a)},{sin(\a)},0)--({cos(\a)},{sin(\a)},\h);

    % Arc angle
    \draw[color=orange,line width=1.1pt]
        plot[domain=0:\a,samples=35]
        ({0.42*cos(\x)},{0.42*sin(\x)},0);

    % Points
    \fill[color=muted] (1,0,\h) circle (1.7pt);
    \fill[color=accent] ({cos(\a)},{sin(\a)},\h) circle (1.9pt);

    % Labels
    \node[color=muted,anchor=west,inner sep=3pt]
        at (0.60,-0.08,\h) {$\vec v$};
    \node[color=accent,anchor=west,inner sep=3pt]
        at ({cos(\a)+0.06},{sin(\a)+0.03},{\h+0.02})
        {$R_z(\theta)\vec v$};
    \node[color=orange,anchor=west,inner sep=2pt]
        at (0.20,0.12,0.03) {$\theta$};
    \node[color=green,anchor=west,inner sep=2pt]
        at (0.05,0.04,0.62) {axe $z$};

\end{tikzpicture}

\end{center}

La rotation se déroule dans le plan $xy$, tandis que l'axe $z$
reste fixe.

% ============================================================
%                         ROTATION X
% ============================================================

\sectiontitle{Rotation autour de l'axe x}

\[
\boxed{
R_x(\theta)=
\begin{pmatrix}
1&0&0\\
0&\cos\theta&-\sin\theta\\
0&\sin\theta&\cos\theta
\end{pmatrix}}
\]

\begin{center}

\tdplotsetmaincoords{68}{122}

\begin{tikzpicture}[tdplot_main_coords,scale=2.45]

    \def\a{52}
    \def\xx{0.9}

    % Axes
    \draw[-{Latex[length=2.5mm]},line width=0.8pt,color=muted]
        (-0.1,0,0)--(1.55,0,0) node[right] {$x$};
    \draw[-{Latex[length=2.5mm]},line width=0.8pt,color=muted]
        (0,-1.25,0)--(0,1.45,0) node[above] {$y$};
    \draw[-{Latex[length=2.5mm]},line width=0.8pt,color=muted]
        (0,0,-0.1)--(0,0,1.5) node[above] {$z$};

    % Axe x
    \draw[color=green,line width=1.1pt]
        (0,0,0)--(1.35,0,0);

    % Vecteurs
    \draw[-{Latex[length=2.7mm]},color=muted,line width=1.3pt]
        (0,0,0)--(\xx,1,0);

    \draw[-{Latex[length=2.7mm]},color=accent,line width=1.35pt]
        (0,0,0)--(\xx,{cos(\a)},{sin(\a)});

    % Construction
    \draw[dashed,color=muted!75,line width=0.7pt]
        (\xx,1,0)--(\xx,0,0);
    \draw[dashed,color=muted!75,line width=0.7pt]
        (\xx,{cos(\a)},{sin(\a)})--(\xx,0,{sin(\a)});

    % Arc angle
    \draw[color=orange,line width=1.1pt]
        plot[domain=0:\a,samples=35]
        (\xx,{0.42*cos(\x)},{0.42*sin(\x)});

    % Points
    \fill[color=muted] (\xx,1,0) circle (1.7pt);
    \fill[color=accent] (\xx,{cos(\a)},{sin(\a)}) circle (1.9pt);

    % Labels
    \node[color=muted,anchor=west,inner sep=3pt]
        at (\xx+0.04,0.55,0.05) {$\vec v$};
    \node[color=accent,anchor=west,inner sep=3pt]
        at (\xx+0.05,{cos(\a)},{sin(\a)+0.03})
        {$R_x(\theta)\vec v$};
    \node[color=orange,anchor=west,inner sep=2pt]
        at (\xx,0.16,0.16) {$\theta$};
    \node[color=green,anchor=west,inner sep=2pt]
        at (0.65,0.04,0.03) {axe $x$};

\end{tikzpicture}

\end{center}

La transformation agit dans le plan $yz$ et laisse $x$ inchangé.

% ============================================================
%                         ROTATION Y
% ============================================================

\sectiontitle{Rotation autour de l'axe y}

\[
\boxed{
R_y(\theta)=
\begin{pmatrix}
\cos\theta&0&\sin\theta\\
0&1&0\\
-\sin\theta&0&\cos\theta
\end{pmatrix}}
\]

\begin{center}

\tdplotsetmaincoords{68}{122}

\begin{tikzpicture}[tdplot_main_coords,scale=2.45]

    \def\a{50}
    \def\yy{0.7}

    % Axes
    \draw[-{Latex[length=2.5mm]},line width=0.8pt,color=muted]
        (-0.15,0,0)--(1.55,0,0) node[right] {$x$};
    \draw[-{Latex[length=2.5mm]},line width=0.8pt,color=muted]
        (0,-1.25,0)--(0,1.45,0) node[above] {$y$};
    \draw[-{Latex[length=2.5mm]},line width=0.8pt,color=muted]
        (0,0,-0.1)--(0,0,1.5) node[above] {$z$};

    % Axe y
    \draw[color=green,line width=1.1pt]
        (0,0,0)--(0,1.25,0);

    % Vecteurs
    \draw[-{Latex[length=2.7mm]},color=muted,line width=1.3pt]
        (0,0,0)--(0,\yy,1);

    \draw[-{Latex[length=2.7mm]},color=accent,line width=1.35pt]
        (0,0,0)--({sin(\a)},\yy,{cos(\a)});

    % Arc angle
    \draw[color=orange,line width=1.1pt]
        plot[domain=0:\a,samples=35]
        ({0.42*sin(\x)},\yy,{0.42*cos(\x)});

    % Points
    \fill[color=muted] (0,\yy,1) circle (1.7pt);
    \fill[color=accent] ({sin(\a)},\yy,{cos(\a)}) circle (1.9pt);

    % Labels
    \node[color=muted,anchor=west,inner sep=3pt]
        at (0.02,\yy,0.62) {$\vec v$};
    \node[color=accent,anchor=west,inner sep=3pt]
        at ({sin(\a)+0.05},\yy,{cos(\a)+0.03})
        {$R_y(\theta)\vec v$};
    \node[color=orange,anchor=west,inner sep=2pt]
        at (0.12,\yy,0.28) {$\theta$};
    \node[color=green,anchor=west,inner sep=2pt]
        at (0.03,0.72,0.03) {axe $y$};

\end{tikzpicture}

\end{center}

Ici, la rotation se déroule dans le plan $xz$, tandis que $y$
reste inchangé.

% ============================================================
%                         OBJET 3D
% ============================================================

\newpage

\sectiontitle{Transformer un objet 3D}

Un objet peut être décrit par ses sommets :

\[
P_1,P_2,\ldots,P_n.
\]

Une matrice $A$ est appliquée à chacun d'eux :

\[
\boxed{P_i'=AP_i}
\]

La transformation de l'objet entier est donc construite à partir
de la transformation de chacun de ses sommets.

\begin{center}

\tdplotsetmaincoords{67}{122}

\begin{tikzpicture}[tdplot_main_coords,scale=2.0]

    % Cube original
    \coordinate (A1) at (-1.45,0,0);
    \coordinate (B1) at (-0.45,0,0);
    \coordinate (C1) at (-0.45,1,0);
    \coordinate (D1) at (-1.45,1,0);
    \coordinate (E1) at (-1.45,0,1);
    \coordinate (F1) at (-0.45,0,1);
    \coordinate (G1) at (-0.45,1,1);
    \coordinate (H1) at (-1.45,1,1);

    \draw[color=muted!55,line width=0.85pt]
        (A1)--(B1)--(C1)--(D1)--cycle;
    \draw[color=muted!55,line width=0.85pt]
        (E1)--(F1)--(G1)--(H1)--cycle;
    \draw[color=muted!45,line width=0.8pt]
        (A1)--(E1) (B1)--(F1) (C1)--(G1) (D1)--(H1);

    % Flèche
    \draw[-{Latex[length=2.8mm]},color=orange,line width=1.35pt]
        (-0.25,0.55,1.3)--(0.35,0.55,1.3);

    \node[color=orange,anchor=south,inner sep=3pt]
        at (0.05,0.55,1.33) {$A$};

    % Cube transformé
    \coordinate (A2) at (0.75,0.18,0.05);
    \coordinate (B2) at (1.72,0.48,0.05);
    \coordinate (C2) at (1.45,1.35,0.05);
    \coordinate (D2) at (0.48,1.05,0.05);
    \coordinate (E2) at (0.75,0.18,1.05);
    \coordinate (F2) at (1.72,0.48,1.05);
    \coordinate (G2) at (1.45,1.35,1.05);
    \coordinate (H2) at (0.48,1.05,1.05);

    \draw[color=accent!65,line width=0.9pt]
        (A2)--(B2)--(C2)--(D2)--cycle;
    \draw[color=accent!65,line width=0.9pt]
        (E2)--(F2)--(G2)--(H2)--cycle;
    \draw[color=accent!55,line width=0.8pt]
        (A2)--(E2) (B2)--(F2) (C2)--(G2) (D2)--(H2);

    \node[color=muted,anchor=north,inner sep=3pt]
        at (-0.95,0.5,0) {objet initial};
    \node[color=accent,anchor=north,inner sep=3pt]
        at (1.1,0.78,0) {objet transformé};

\end{tikzpicture}

\end{center}

% ============================================================
%                         COMPOSITION
% ============================================================

\sectiontitle{Composer plusieurs rotations}

Supposons que nous appliquions successivement :

\[
R_x(\alpha),
\qquad
R_y(\beta),
\qquad
R_z(\gamma).
\]

La transformation globale peut être écrite :

\[
\boxed{
A=R_z(\gamma)R_y(\beta)R_x(\alpha)
}
\]

Un vecteur devient alors :

\[
\boxed{\vec v'=A\vec v}
\]

\begin{importantbox}
L'ordre des rotations est important.

\[
\boxed{R_xR_y\neq R_yR_x}
\]
\end{importantbox}

% ============================================================
%                         EXEMPLE
% ============================================================

\sectiontitle{Exemple numérique}

Pour une rotation autour de $z$ de $90^\circ$ :

\[
R_z(90^\circ)=
\begin{pmatrix}
0&-1&0\\
1&0&0\\
0&0&1
\end{pmatrix}.
\]

Prenons :

\[
\vec v=
\begin{pmatrix}
1\\
0\\
2
\end{pmatrix}.
\]

Alors :

\[
\begin{aligned}
\vec v'
&=
\begin{pmatrix}
0&-1&0\\
1&0&0\\
0&0&1
\end{pmatrix}
\begin{pmatrix}
1\\
0\\
2
\end{pmatrix}\\[0.3cm]
&=
\begin{pmatrix}
0\\
1\\
2
\end{pmatrix}.
\end{aligned}
\]

La coordonnée $z$ reste inchangée, tandis que le vecteur tourne
dans le plan $xy$.

% ============================================================
%                         INFORMATIQUE
% ============================================================

\sectiontitle{Lien avec la 3D informatique}

Les matrices de transformation sont utilisées pour manipuler des
objets et des vecteurs dans de nombreux systèmes de calcul 3D.

Une chaîne conceptuelle peut être écrite :

\[
\boxed{
\text{modèle}
\rightarrow
\text{transformation}
\rightarrow
\text{caméra}
\rightarrow
\text{projection}
}
\]

Une scène 3D peut donc être pensée comme un ensemble de points
auxquels on applique des transformations mathématiques.

% ============================================================
%                         CONCLUSION
% ============================================================

\sectiontitle{Conclusion}

Nous sommes partis du repère spatial :

\[
(x,y,z).
\]

Puis nous avons introduit les vecteurs :

\[
\vec v=
\begin{pmatrix}
 x\\
 y\\
 z
\end{pmatrix}.
\]

Une matrice $3\times3$ permet de transformer ces coordonnées :

\[
\boxed{
\begin{pmatrix}
x'\\
y'\\
z'
\end{pmatrix}
=
\begin{pmatrix}
a&b&c\\
d&e&f\\
g&h&i
\end{pmatrix}
\begin{pmatrix}
x\\y\\z
\end{pmatrix}
}
\]

Nous avons ensuite construit les trois rotations fondamentales :

\[
\boxed{
R_x(\theta),
\qquad
R_y(\theta),
\qquad
R_z(\theta)
}
\]

Enfin, plusieurs transformations peuvent être combinées par
multiplication matricielle.

\begin{center}

\vspace{0.8cm}

{\Large\color{accent}
Géométrie $\rightarrow$ vecteurs $\rightarrow$ matrices $3\times3$\par
}

\vspace{0.4cm}

{\color{muted}
Une transformation spatiale peut être décrite par l'algèbre.\par
}

\end{center}

\end{document}
